\subsection{Good filtrations}

Let A be a filtered commutative ring, E a filtered A-module and \((A_{n})\) and \((E_{n})\)

the respective filtrations of A and E; suppose that \(A_{0} = A\) . In the polynomial ring A{[}X{]} the set \(A' = \sum_{n \geqslant 0} A_{n} X^{n}\) is a graded sub-A-algebra of type N; the subgroup \(E' = \sum_{n \geqslant 0} E_{n} \otimes_{A} A X^{n}\) of \(E \otimes_{A} A[X]\) is a graded \(A'\) -module of type N, since

\[
\mathbf {A} _ {m} \mathbf {X} ^ {m} (\mathrm{E} _ {n} \otimes_ {\mathbf {A}} \mathbf {A X} ^ {n}) \subset (\mathbf {A} _ {m} \mathrm{E} _ {n} \otimes_ {\mathbf {A}} \mathbf {A X} ^ {m + n}) \subset \mathrm{E} _ {m + n} \otimes_ {\mathbf {A}} \mathbf {A X} ^ {m + n}.
\]

DEFINITION 1. Let A be a commutative ring, m an ideal of A, E an A-module and \((\mathrm{E}_{n})\) a filtration on the additive group E consisting of submodules of E. The filtrations \((\mathrm{E}_{n})\) is called m-good if:

\begin{enumerate}
\def\labelenumi{(\arabic{enumi})}
\item
  \(\mathfrak{m}\mathrm{E}_n\subset \mathrm{E}_{n + 1}\) for all \(n\in \mathbf{Z}\) ;
\item
  there exists an integer \(n_0\) such that \(\mathfrak{m}\mathrm{E}_n = \mathrm{E}_{n+1}\) for \(n \geqslant n_0\) .
\end{enumerate}

Then by induction on \(q\) , \(\mathfrak{m}^q\mathrm{E}_n = \mathrm{E}_{n + q}\) for \(n \geqslant n_0\) , \(q \geqslant 1\) . Note that condition (1) means that the filtration \((\mathrm{E}_n)\) is compatible with the A-module structure on E if A is given the m-adic filtration. Clearly, on every A-module E, the m-adic filtration is m-good. If a filtration on an A-module E is m-good, the quotient filtration on any quotient module of E is m-good.

THEOREM 1. Let A be a commutative ring, m an ideal of A, E an A-module and \((\mathrm{E}_{n})\) a filtration of the additive group E consisting of finitely generated sub-A-modules. Suppose that \(mE_{n} \subset E_{n+1}\) for all n.~Let \(A'\) be the graded subalgebra \(\sum_{n \geqslant 0} m^{n} X^{n}\) of A{[}X{]} and \(E'\) the graded \(A'\) -module \(\sum_{n \geqslant 0} E_{n} \otimes_{A} AX^{n}\) . The two following conditions are equivalent:

\begin{enumerate}
\def\labelenumi{(\alph{enumi})}
\item
  The filtration \((\mathbf{E}_n)\) is m-good.
\item
  E' is a finitely generated A'-module.
\end{enumerate}

Suppose that \(mE_{n-1} = E_n\) for \(n > n_0 \geqslant 0\) . For \(i \leqslant n_0\) , let \((e_{ij})_{1 \leqslant j \leqslant r_i}\) be a finite system of generators of the A-module \(E_i\) . As the A-module \(E_n \otimes_A X^n\) is generated by the elements \(e_{nj} \otimes X^n\) for \(0 \leqslant n \leqslant n_0\) and is equal to

\[
\mathfrak {m} ^ {n - n _ {0}} \mathrm{E} _ {n _ {0}} \otimes_ {\mathbf {A}} \mathrm{AX} ^ {n}
\]

for \(n > n_0\) , the A'-module E' is generated by the elements \(e_{nj} \otimes \mathbf{X}^n\) for \(0 \leqslant n \leqslant n_0\) and \(1 \leqslant j \leqslant r_n\) ; then it is certainly finitely generated.

Conversely, if \(\mathbf{E}'\) is a finitely generated \(\mathbf{A}'\) -module, it is generated by a finite family of elements of the form \(e_k \otimes \mathbf{X}^{n(k)}\) , where \(e_k \in \mathbf{E}_{n(k)}\) . Let \(n_0\) be the greatest of the integers \(n(k)\) . Then for \(n \geqslant n_0\) and \(f \in \mathbf{E}_n\) ,

\[
f \otimes \mathbf {X} ^ {n} = \sum_ {k} t _ {k} (e _ {k} \otimes \mathbf {X} ^ {n (k)})
\]

where \(t_k \in \mathbf{A}'\) ; replacing if need be \(t_k\) by its homogeneous component of degree \(n - n(k)\) , we may assume that \(t_k = a_k \mathbf{X}^{n - n(k)}\) , where \(a_k \in \mathfrak{m}^{n - n(k)}\) .

As the unique element \(X^{n}\) forms a basis of the A-module \(AX^{n}\) , the equation \(f \otimes X^{n} = \left( \sum_{k} a_{k} e_{k} \right) \otimes X^{n}\) implies \(f = \sum_{k} a_{k} e_{k}\) . Then \(E_{n} \subset m^{n-n_{0}} E_{n_{0}}\) ; since the opposite inclusion is obvious,

\[
\mathrm{E} _ {n} = \mathfrak {m} ^ {n - n _ {0}} \mathrm{E} _ {n _ {0}},
\]

whence \(\mathbf{E}_n = m\mathbf{E}_{n - 1}\) for \(n > n_0\)

LEMMA 1. Let A be a commutative Noetherian ring and m an ideal of A. Then the subring \(\mathbf{A}' = \sum_{n\geqslant 0}\mathfrak{m}^n\mathbf{X}^n\) of A{[}X{]} is Noetherian.

A' is an A-algebra generated by mX; as A is Noetherian, mX is a finitely generated A-module and the conclusion then follows from § 2, no. 10, Corollary 3 to Theorem 2.

PROPOSITION 1. Let A be a commutative Noetherian ring and m an ideal of A; let A be given the m-adic filtration. Let E, F be two filtered A-modules and j: F → E an injective homomorphism compatible with the filtrations. If E is finitely generated and its filtration is m-good, then F is finitely generated and its filtration is m-good.

As F is isomorphic to a submodule of E, it is finitely generated since A is Noetherian and E is finitely generated. Let \((\mathbf{E}_{n})\) , \((\mathbf{F}_{n})\) be the respective filtrations on E and F, which consist of finitely generated submodules; preserving the notation of Lemma 1, we set \(E' = \sum_{n \geqslant 0} E_{n} \otimes_{A} AX^{n}\) , \(F' = \sum_{n \geqslant 0} F_{n} \otimes_{A} AX^{n}\) ; as by hypothesis \(F_{n}\) is isomorphic to a submodule of \(E_{n}\) , we see that \(F'\) is isomorphic to a submodule of \(E'\) . By Theorem 1, \(E'\) is a finitely generated \(A'\) -module and hence so is \(F'\) since \(A'\) is Noetherian (Lemma 1). Hence the conclusion by virtue of Theorem 1.

COROLLARY 1 (the Artin-Rees Lemma). Let A be a commutative Noetherian ring, m an ideal of A, E a finitely generated A-module and F a submodule of E. The filtration induced on F by the m-adic filtration on E is m-good.

In other words, there exists an integer \(n_{0}\) such that

\[
\mathfrak {m} ((\mathfrak {m} ^ {n} \mathrm{E}) \cap \mathrm{F}) = (\mathfrak {m} ^ {n + 1} \mathrm{E}) \cap \mathrm{F}\tag{1}
\]

for all \(n \geqslant n_0\) .

COROLLARY 2. Let A be a commutative Noetherian ring and a, b two ideals of A. There exists an integer h \textgreater{} 0 such that \(a^{h} \cap b \subset ab\) .

There exists \(n\) such that \(\mathfrak{a}^{n+1} \cap \mathfrak{b} = \mathfrak{a}(\mathfrak{a}^n \cap \mathfrak{b}) \subset \mathfrak{ab}\) by Corollary 1 applied to \(\mathbf{E} = \mathbf{A}, \mathbf{F} = \mathfrak{b}\) .

COROLLARY 3. Let A be a commutative Noetherian ring, m an ideal of A and x an element of A which is not a divisor of 0. There exists an integer \(k > 0\) such that, for all \(n \geqslant k\) , the relation \(xy \in \mathfrak{m}^n\) implies \(y \in \mathfrak{m}^{n - k}\) .

Corollary 1 applied to \(\mathbf{E} = \mathbf{A}\) , \(\mathbf{F} = \mathbf{A}x\) shows that there exists \(k\) such that, for all \(n \geqslant k\) , \(\mathfrak{m}^n \cap \mathbf{A}x = \mathfrak{m}^{n - k}(\mathfrak{m}^k \cap \mathbf{A}x)\) . Then, if \(xy \in \mathfrak{m}^n\) ,

\[
x y \in \mathfrak {m} ^ {n} \cap A x \subset \mathfrak {m} ^ {n - k} x
\]

and, as \(x\) is not a divisor of 0, we deduce that \(y \in \mathfrak{m}^{n - k}\) .

In the notation of transporters (Chapter I, § 2, no. 10), the conclusion of Corollary 3 reads

\[
\mathfrak {m} ^ {n}: \mathrm{A} x \subset \mathfrak {m} ^ {n - k}.\tag{2}
\]

COROLLARY 4. Let A be a commutative Noetherian ring, m an ideal of A, E a finitely generated A-module and \((\mathbf{E}_{n})\) and \((\mathbf{E}_{n}^{\prime})\) two filtrations consisting of submodules of E. Suppose that the filtrations \((\mathbf{E}_{n})\) and \((\mathbf{E}_{n}^{\prime})\) are compatible with the A-module structure on E when A is given the m-adic filtration. If the filtration \((\mathbf{E}_{n})\) is m-good and \(E_{n}^{\prime} \subset E_{n}\) for all \(n \in Z\) , the filtration \((\mathbf{E}_{n}^{\prime})\) is m-good. This a special case of Proposition 1.

LEMMA 2. Let A, B be two commutative Noetherian rings, \(\phi: \mathrm{A} \to \mathrm{B}\) a ring homomorphism, E a finitely generated A-module and F a finitely generated B-module. Then \(\operatorname{Hom}_{\mathbf{A}}(\mathrm{E}, \phi_*(\mathrm{F}))\) is a finitely generated B-module.

There exists by hypothesis a surjective A-homomorphism \(v: \mathbf{A}^n \to \mathbf{E}\) ; the mapping \(u \mapsto u \circ v\) of \(\operatorname{Hom}_{\mathbf{A}}(\mathbf{E}, \phi_*(\mathbf{F}))\) to \(\operatorname{Hom}_{\mathbf{A}}(\mathbf{A}^n, \phi_*(\mathbf{F}))\) is therefore injective and, as B is Noetherian, it is sufficient to prove that \(\operatorname{Hom}_{\mathbf{A}}(\mathbf{A}^n, \phi_*(\mathbf{F}))\) is a finitely generated B-module; which is immediate since it is isomorphic to \(\mathbf{F}^n\) .

PROPOSITION 2. Let A be a commutative Noetherian ring, m an ideal of A and E, F two finitely generated A-modules. If \((\mathbf{F}_{n})\) is an m-good filtration on F, the submodules \(\operatorname{Hom}_{\mathbf{A}}(\mathbf{E}, \mathbf{F}_{n})\) form an m-good filtration on the A-module \(\operatorname{Hom}_{\mathbf{A}}(\mathbf{E}, \mathbf{F})\) . As \(m^{k}F_{n}\subset F_{n+k}\) for \(n\in Z, k\geqslant0\) , it is also true that

\[
\mathfrak {m} ^ {k} \operatorname{Hom} _ {\mathbf {A}} (\mathrm{E}, \mathrm{F} _ {n}) \subset \operatorname{Hom} _ {\mathbf {A}} (\mathrm{E}, \mathrm{F} _ {n + k});
\]

the family \((\mathrm{Hom}_{\mathbf{A}}(\mathrm{E},\mathrm{F}_{n}))_{n\in\mathbf{Z}}\) is then a filtration on \(\mathrm{Hom}_{\mathbf{A}}(\mathrm{E},\mathbf{F})\) compatible with its module structure over the ring A filtered by the m-adic filtration. Since E is finitely generated, there exists an integer r\textgreater0 and a surjective A-homomorphism \(u:A^{r}\rightarrow E\) which defines an injective A-homomorphism

\[
v = \operatorname{Hom} (u, 1 _ {\mathrm{F}}): \operatorname{Hom} _ {\mathrm{A}} (\mathrm{E}, \mathrm{F}) \rightarrow \operatorname{Hom} _ {\mathrm{A}} (\mathrm{A} ^ {r}, \mathrm{F});
\]

clearly \(v\) is compatible with the filtrations \((\mathrm{Hom}_{\mathbf{A}}(\mathbf{E},\mathbf{F}_n))\) and \((\mathrm{Hom}_{\mathbf{A}}(\mathbf{A}^r,\mathbf{F}_n))\) . As \(\mathrm{Hom}_{\mathbf{A}}(\mathbf{E},\mathbf{F})\) and \(\mathrm{Hom}_{\mathbf{A}}(\mathbf{A}^r,\mathbf{F})\) are finitely generated (Lemma 2), it is sufficient by virtue of Proposition 1 to show that the filtration \((\mathrm{Hom}_{\mathbf{A}}(\mathbf{A}^{r},\mathbf{F}_{n}))\) is m-good; but this is immediate by virtue of the existence of the canonical isomorphism \(\mathrm{Hom}_{\mathbf{A}}(\mathbf{A}^{r},\mathbf{F}_{n})\to \mathbf{F}_{n}^{r}\) and the fact that the relation \(\mathfrak{mF}_n = \mathbf{F}_{n + 1}\) implies \(\mathfrak{m}(\mathbf{F}_n^r) = (\mathfrak{m}\mathbf{F}_n)^r = \mathbf{F}_{n + 1}^r\) (Algebra, Chapter II, § 3, no. 7, Remark).

PROPOSITION 3. Let A be a Noetherian ring and m an ideal of A such that A is Hausdorff and complete with respect to the m-adic topology. Let E be a filtered A-module over the filtered ring A, the filtration ( \(E_{n}\) ) of E being such that \(E_{0} = E\) and E is Hausdorff with respect to the topology defined by ( \(E_{n}\) ). Then the following conditions are equivalent:

\begin{enumerate}
\def\labelenumi{(\alph{enumi})}
\item
  E is a finitely generated A-module and \((\mathbf{E}_{n})\) is an m-good filtration.
\item
  gr(E) is a finitely generated gr(A)-module.
\item
  For all \(n \geqslant 0\) , \(\mathrm{gr}_n(\mathbf{E})\) is a finitely generated A-module and there exists \(n_0\) such that for \(n \geqslant n_0\) the canonical homomorphism

\[
\operatorname{gr} _ {1} (\mathrm{A}) \otimes_ {\mathrm{A}} \operatorname{gr} _ {n} (\mathrm{E}) \rightarrow \operatorname{gr} _ {n + 1} (\mathrm{E})\tag{3}
\]

is surjective.
\end{enumerate}

It follows immediately from the definitions that (a) implies (c). The fact that (b) implies (c) is a consequence of § 1, no. 3, Lemma 1; conversely, if (c) holds, clearly gr(E) is generated as a gr(A)-module by the sum of the gr \(_{p}\) (E) for p ≤ n \(_{0}\) and hence by hypothesis admits a finite system of generators. It remains to prove that (c) implies (a); as the gr \(_{n}\) (E) are finitely generated and E \(_{0}\) = E, clearly first, by induction on n, E/E \(_{n}\) is a finitely generated A-module for all n; it will therefore be sufficient to prove that, for n \textgreater{} n \(_{0}\) , E \(_{n}\) is a finitely generated A-module and that mE \(_{n}\) = E \(_{n+1}\) . Now, consider the A-module E \(_{n+1}\) with the exhaustive and separated filtration (E \(_{n+k}\) )(k ≥ 1); mE \(_{n}\) ⊂ E \(_{n+1}\) ; hypothesis (c) implies that the image of mE \(_{n}\) in gr \(_{n+1}\) (E) = E \(_{n+1}\) /E \(_{n+2}\) is equal to gr \(_{n+1}\) (E) and generates the graded gr(A)-module gr(E \(_{n+1}\) ). As gr \(_{n+1}\) (E) is by hypothesis a finitely generated A-module, it follows from § 2, no. 9, Proposition 12 that mE \(_{n}\) = E \(_{n+1}\) and that E \(_{n+1}\) is a finitely generated A-module.

